\section*{Concluding remarks}

This Supplementary Information was written so that every claim of the main text can be traced to its source and checked independently. Supplementary Tables S1 and S2 give the subsets selected for each benchmark problem and the notation of the method, and Supplementary Note~1 shows that our simulation of the measured outputs agrees with the reference simulations of the benchmark. These materials connect the results reported in the paper with the data and code from which they were obtained.

The remaining material documents the formal verification. Supplementary Table~S3 lists the 31 modules of the Lean~4 formalisation in the order in which they depend on each other, and Supplementary Note~2 follows that order to explain, in mathematical language, how the 289 theorems and lemmas build up the guarantees of the method: from the energy bound on the error of the reduced model and the equivalence between $\cos\Delta$ and the best relative fitting error, through the differentiability of ODE solutions with respect to the parameters, to certificates that can be computed for ODE models and to the hypotheses of the benchmark models. Supplementary Note~3 and Table~S4 give the statement and meaning of every theorem, so that a reader can compare each formal statement with the mathematics it encodes without reading Lean code. Supplementary Note~4 and Table~S5 show, model by model, the data on which the Lean kernel executed its checks and the theorem that results.

Taken together, the formalisation establishes that the steps of the method are mathematically sound and that 22 of the benchmark models satisfy their hypotheses; the numerical values reported in the paper are computed by the accompanying Python code, and the formal results state precisely which conclusions follow from them. Everything described here can be reproduced from the repository: the proofs compile with a single command, the supplementary tables and notes are generated from the Lean sources and the result files, and the analyses can be rerun from the published benchmark. We hope that this combination of a public benchmark, open code and machine-checked mathematics makes the reduction method easy to verify, to reuse and to extend to other dynamical models.
